\documentclass[11pt,a4paper]{article}
\usepackage{../pelm}
\begin{document}
\paper{Final Examination}{Whole course, weighted towards Weeks 6--13 \quad$\cdot$\quad 40 questions $\times$ 2.5 points = 100 points, 60\% of the course grade}{40}{90 minutes}

\begin{questions}

% ============================================== Part 1: mechanism
\q{Which best describes what a large language model does?}
  {Retrieves the most relevant stored fact}
  {Produces a probability for every possible next token, then one is selected}
  {Applies rules written by its developers}
  {Searches the internet and summarises the results}

\q{Why is hallucination not simply a bug awaiting a fix?}
  {The models are still too small}
  {There is no step in the operation at which truth is checked; plausibility is what is optimised}
  {Training data has not yet been cleaned}
  {Users write insufficiently precise prompts}

\q{In which situation is hallucination \textbf{most} likely?}
  {A widely documented topic in English}
  {A specific citation on a thinly covered topic}
  {A request for a summary of provided text}
  {A request to rewrite a paragraph more simply}

\q{A question is phrased ``Explain why X causes Y'', when in fact X does not cause Y. What is the
   likely result?}
  {The model corrects the premise}
  {The model produces an explanation, accepting the premise}
  {The model refuses to answer}
  {The model reports low confidence}

\q{Why does the same content cost more tokens in Turkish than in English?}
  {Turkish words are longer on average}
  {Agglutinative stacking produces assembled words that appeared rarely in training}
  {Turkish requires a separate tokenizer pass}
  {Turkish uses more punctuation}

\q{An assistant is asked to count letters in a word and gets it wrong. Why?}
  {The arithmetic exceeded its capability}
  {It never perceived individual letters --- it operates on tokens}
  {The context window overflowed}
  {The temperature was set too high}

\q{Which is \textbf{not} held in the context window?}
  {The conversation so far}
  {Attached documents}
  {Hidden instructions added by the product}
  {The model's training data}

\q{What does it mean that these systems are not deterministic?}
  {They cannot perform arithmetic}
  {The same prompt may produce different answers on different runs}
  {They cannot be tested}
  {Their output depends on the time of day}

\q{Why is repeating a question and receiving a similar answer weak evidence of correctness?}
  {The model remembers its previous answer}
  {Consistency measures how strongly a pattern is represented, not whether it is true}
  {The second answer is always shorter}
  {The context window has changed}

\q{Which best states the relationship between plausibility and truth in these systems?}
  {They are equivalent}
  {They correlate, because training text was mostly written by people trying to be accurate}
  {They are unrelated}
  {Truth is checked and plausibility is a by-product}

\q{Adding ``be accurate and use only reliable sources'' to a prompt about a thinly covered topic
   will most likely:}
  {Improve factual accuracy}
  {Produce a refusal}
  {Make the answer sound more authoritative without improving it}
  {Trigger a web search}

\q{Which of the four levers should be used when the output has the right content in the wrong
   shape?}
  {Context}
  {Task}
  {Format}
  {Examples}

\q{Which is a well-formed negative constraint?}
  {Do not be vague}
  {Do not write anything unprofessional}
  {Do not cite any case or regulation by name}
  {Do not disappoint the reader}

% ============================================== Part 2: diagnosis
\q{An answer about Turkish procedure describes, accurately, the procedure of another country.
   This is:}
  {Fabrication}
  {Displacement}
  {Omission}
  {Silent failure}

\q{A legal memo contains a case name in correct citation format for a case that was never
   decided. This is:}
  {Fabrication}
  {Displacement}
  {Over-reading}
  {Indirect injection}

\q{Which failure is hardest to detect outside your own area of expertise, and why?}
  {Fabrication --- because citations look real}
  {Displacement --- because foreign systems are unfamiliar}
  {Omission --- because there is no incorrect statement to identify}
  {Overconfidence --- because tone is uniform}

\q{A grounded assistant answers a question your documents do not cover, producing a confident
   response with no quotation. This is:}
  {Retrieved the wrong passage}
  {Over-read the passage}
  {Answered from memory}
  {Silent failure}

\q{Your source says ``the committee \emph{may} require additional evidence''. The assistant
   reports ``the committee \emph{requires} additional evidence''. This is:}
  {Fabrication}
  {Over-reading the passage}
  {Answered from memory}
  {Displacement}

\q{A grounded answer quotes a real sentence from your document, but the sentence comes from an
   unrelated section and does not support the claim. This is:}
  {Fabrication}
  {Retrieved the wrong passage}
  {Omission}
  {Leakage}

\q{A three-step workflow receives faulty input at step one and produces a complete, confident,
   well-formatted output. Nothing reports an error. This is:}
  {Wrong stopping point}
  {Silent failure}
  {Context overflow}
  {Direct injection}

\q{What do all nine failure modes covered in this course have in common?}
  {They occur only in long conversations}
  {They are invisible in the output}
  {They can be prevented by better prompting}
  {They occur only in grounded systems}

\q{A generated browser tool produces a subtly wrong number, consistently, behind a clean
   interface. Why is consistency dangerous here?}
  {It makes the error harder to reproduce}
  {It makes the tool appear reliable}
  {It increases the token cost}
  {It causes the tool to fail eventually}

\q{Which check would most reliably expose a poisoned document summary?}
  {Asking the model whether it was manipulated}
  {Regenerating the summary several times}
  {Checking each claim against a quoted passage in the source}
  {Reducing the temperature}

% ============================================== Part 3: design and judgment
\q{What is the most significant thing lost when a chatbot is replaced by an agent?}
  {Response speed}
  {The unconscious checking you performed on every intermediate output}
  {The ability to attach documents}
  {Access to the context window}

\q{If each step is 95\% reliable, approximately what is end-to-end reliability over ten steps?}
  {95\%}
  {80\%}
  {60\%}
  {40\%}

\q{Why is compounding error worse than independent random error?}
  {Errors are reported only at the end}
  {Each step consumes the previous step's output, so early errors are inherited}
  {Agents use more tokens per step}
  {Later steps are inherently less reliable}

\q{Which control does \textbf{not} rely on the system behaving as intended?}
  {Instructing it to be careful}
  {Asking it to verify its own output}
  {Connecting only the narrowest tools required}
  {Requesting quotations for each claim}

\q{Where should a human be placed in an AI workflow?}
  {At the end, reviewing the output}
  {Before the irreversible action}
  {At the start, writing the prompt}
  {Wherever capacity allows}

\q{Which is \textbf{not} a criterion for a good first automation candidate?}
  {The task is genuinely repetitive}
  {The answers exist in documents you can point at}
  {A wrong answer is recoverable}
  {The task is the most intellectually difficult in your field}

\q{Why should the harmed person be named specifically rather than as a category?}
  {Institutional policy requires it}
  {Specificity changes the design; abstraction does not}
  {It shortens the risk assessment}
  {Categories cannot be defined precisely}

\q{In a well-designed AI workflow, most decisions concern:}
  {Which model to use}
  {What the system must not do}
  {How to phrase the prompt}
  {Which tool provider to select}

\q{A Domain Pack improves all of the following \textbf{except}:}
  {Local vocabulary}
  {Recognisable structure}
  {Conventional restraint}
  {Factual accuracy}

\q{Why can a Domain Pack increase risk as well as usefulness?}
  {It consumes context window space}
  {Professional formatting makes wrong content more persuasive}
  {It slows the system down}
  {It prevents the system from refusing}

\q{Why do professional structures such as SOAP and IRAC help when working with these systems?}
  {They are shorter than free prose}
  {Each makes a particular kind of error visible}
  {They are familiar to the model}
  {They reduce token usage}

\q{What is the root cause of prompt injection?}
  {Insufficient safety fine-tuning}
  {The model receives one stream of tokens and cannot reliably separate instructions from data}
  {Context windows are too small}
  {Users upload untrusted files}

\q{Why is indirect injection more serious for an ordinary user than direct injection?}
  {It is technically harder to perform}
  {The malicious instruction is hidden in content the user never sees, so there is no point at which they could refuse}
  {It cannot be defended against at all}
  {It only affects grounded systems}

\q{Which attacker goal requires an agent rather than a chatbot?}
  {Changing the content of a summary}
  {Extracting other material from the context}
  {Causing an action such as sending or deleting}
  {Revealing the system instructions}

\q{Which statement about anonymisation is correct?}
  {Removing names is sufficient}
  {A rare condition, a date and a small department can identify someone even with names removed}
  {Anonymised data may be uploaded to any service}
  {Anonymisation is required only for research}

\q{Under this course's account of academic integrity, the central question is:}
  {Whether AI was used at all}
  {What the assessment is measuring, and what you are certifying by submitting}
  {Whether the use can be detected}
  {How much of the text was generated}

\q{What is the strongest argument for occasionally drafting work yourself?}
  {It is faster for short documents}
  {The judgment that lets you evaluate generated output is built by doing the work}
  {Institutional rules require it}
  {Generated drafts are usually poor}

\end{questions}
\end{document}
